\chapter{Introduction}
\pagenumbering{arabic}

\section{overview}
In this thesis we attempt to extend the {\em reflection/transmission\/} 
matrix method of {\em Kennett\/}~(1974) to laterally heterogeneous media. 
We assume a $2-$D model for the material heterogeneity of the medium, but 
the wavefield is allowed to spread out in $3-$D, which is
sometimes referred to as a $2 \slantfrac{1}{2}-$D model. We assume that 
the heterogeneity is an extra level of structure that overlays an otherwise
horizontally stratified half space. We call the stratified component of 
the model, the {\em reference model\/}, and the remaining variations of 
material properties are relegated to the {\em residual model\/}. Wave 
propagation in the reference model (which will be called the {\em 
reference field\/}) is now well understood [e.g. {\em Aki \& 
Richards\/}~(1980); {\em Kennett\/}~(1983)].

The residual structure, {\em scatters\/} the reference wavefield into the 
{\em actual wavefield\/}, which is the quantity of interest. The 
{\em difference between the actual wavefield and the reference wavefield 
is the scattered field\/}.  This scattered field may be imagined to 
derive from  a {\em secondary\/} or {\em virtual\/} source distribution.
Consistent with {\em Huyghens' principle\/}, we assume
that the secondary source distribution, may 
be discretized, as a {\em regular lattice of radiating multipoles\/}.
 Due 
to the {\em diffuse\/} nature of seismic scattering, the secondary source 
distribution {\em covers\/} the spatial support of the residual structure.
Unfortunately, the secondary source field is unknown, 
because it depends  on the actual wavefield, which is yet to be determined. 
In this thesis we propose to determine the secondary source field,  by
summing the successive interactions of the incident field with 
heterogeneity. We refer to this procedure as: {\em Summing the multiple 
scattering series\/}.

This program has its origins in the quantum mechanics formulation
[e.g. {\em Agassi \& Gal\/}~\shortcite{AgassiandGal},
{\em Keister\/}~\shortcite{Keister}, {\em Gonis\/}~\shortcite{Gonis}],
of scalar wave scattering,
and the elastic wave equation extension has been investigated in {\em 
Kennett\/}~(1986), where it was shown how to calculate the {\em single 
scattering\/} approximation of the actual wavefield. In {\em 
Kennett\/}~(1986), the superimposed heterogeneity was subdivided into 
{\em thin slabs\/}. In the single scattering approximation to a 
reflection seismogram, the reflected field is comprised of only the {\em 
primary backscatter interactions\/} of the incident reference field with 
the thin slabs of superimposed 
heterogeneity, while incorporating all transmission and reflections processes 
in the reference model.

The main focus of this thesis, is the investigation of:
\begin{itemize}
\item  the efficient computation of  multiple scattering interactions of the 
incident wavefield with the residual structure
\item  methods of summing the multiple scattering partial sums  in order to secure 
rapid convergence of the multiple scattering series
\end{itemize}
A choice must be made concerning the 
representation of the superimposed heterogeneity (i.e. the residual 
model). The most obvious choice is to divide the residual model into thin 
slabs, as employed in  {\em Kennett\/}~(1986), and approximate the scattering 
properties of each slab  using the {\em Born\/} or {\em single 
scattering\/} approximation. This model is based on the assumption that 
there is more energy involved in {\em inter-slab\/} scattering interactions than 
{\em intra-slab\/} scattering interactions. This representation ignores 
the possibility of {\em side scattering\/} between two point within the 
same thin layer.

%%%%%%%%%%%%%%%%%%%%%%%%%%%%%%%%%%%%%%%%%%%%%%%%%%%%%%%%%%%%%%%%%%%%%%%%%%%%%%%%%
%%%%%%%%%%%%%%%%%%%%%%%%%%%%%%%%%%%%%%%%%%%%%%%%%%%%%%%%%%%%%%%%%%%%%%%%%%%%%%%%%%%%
\begin{figure}
\HideDisplacementBoxes
%\ShowDisplacementBoxes
\centerline{\BoxedEPSF{Rdpp.EPSF scaled 1000}}
\caption{Multiple scattering Rpp for $6\times 6$ point scatterers}
\protect\label{Rdpp}
\end{figure}
\begin{figure}
\HideDisplacementBoxes
%\ShowDisplacementBoxes
\centerline{\BoxedEPSF{dRdpp.EPSF scaled 1000}}
\caption{The side scattering component of Rpp for $6\times6$ point scatterers}
\protect\label{dRdpp}
\end{figure}
%%%%%%%%%%%%%%%%%%%%%%%%%%%%%%%%%%%%%%%%%%%%%%%%%%%%%%%%%%%%%%%%%%%%%%%%%%%%%%%%%%%%
%%%%%%%%%%%%%%%%%%%%%%%%%%%%%%%%%%%%%%%%%%%%%%%%%%%%%%%%%%%%%%%%%%%%%%%%%%%%%%%%%

In Fig.\ref{Rdpp} we depict the in-plane, downward incident $P-$wave to 
upward reflected $P-$wave response for $36$ point lattice of scatterers arranged in a square, 
by summing the multiple scattering series to a convergence, for $101$ 
incident plane waves of varying horizontal wavenumber. In Fig.\ref{dRdpp} 
we depict the contribution in the experiment of Fig.\ref{Rdpp} due to 
side scattering between point scatterers in the same row. Certainly, for a 
low order approximation, we could neglect the side scatter component 
depicted by Fig.\ref{dRdpp}
altogether, and we will see that there are significant computational 
advantages  in doing just that. However, if we require a more accurate 
calculation of the the true response, we will need to incorporate these 
side scattering components.  Unfortunately, there 
is a problem with this approach. 
If we employ conventional plane wave techniques, the 
Green's tensor which is needed to model the side scatter interactions 
diverges for lattice points within the same horizontal row, 
{\em especially in the limiting case where the horizontal offset between 
a pair of points tends to zero, so that the scattering point is interacting with itself\/}.

We show how to model side scattering, and particularly, 
self interaction, where the Green's tensor needed to transport the wave 
field along its scattering path is {\em strongly singular\/}. This 
involves finding new representations of the Green's tensor, and 
evaluating the self interaction terms when the source point and field 
point of the Green's tensor coincide. We have chosen to 
model the residual medium as a {\em tessellation\/} of small squares with 
uniform material properties, rather than thin slabs. Consistent with the 
Rayleigh approximation (Wu \& Aki, 1985),   the secondary wavefield, 
is modeled as a multipolar point force, situated in 
the middle of the respective squares.  We compute the secondary source 
field (i.e. the multipole amplitudes) using the {\em pseudo-spectral method\/}.
Indeed, our multiple scattering method shares some of the characteristics 
of the {\em elastic wave pseudo-spectral method\/} ({Kosloff, Reshef \& 
Loewenthal}, 1984), with both methods requiring similar amounts of FFT 
computation. 

Generally, the  pseudo-spectral method requires less 
computation than our {\em multiple scattering method\/} for elastic media with strong superimposed 
heterogeneity (say $10 \%$ contrast relative to reference model), because 
it uses {\em second order time stepping\/}, which is very efficient to 
apply numerically. In the multiple scattering method of this thesis, 
the secondary source field is computed in {\em pseudo-spectral\/} fashion, 
but instead of time stepping, we have to compute a {\em monochromatic\/} 
multi-source/receiver seismogram, which is considerably more expensive 
than the analogous time stepping operation of the pseudo-spectral method.
However, there are generally more time steps in the pseudo-spectral 
method than there are complex frequencies in the multiple scattering 
method.
 

When using second order time stepping, we must ensure that numerical errors 
are not amplified from one time step to 
the next. The time step ${\mit \Delta }t$ is constrained by the numerical 
stability requirement that 
${\mit \Delta}t \le 0.26 {\mit \Delta }{ x}_{min}/{c}_{max}$  (Daudt, 
Braile,Nowack, \& Chiang, 1989),  where ${\mit \Delta }{ x}_{min}$ is the 
minimum spatial grid spacing and ${c}_{max}$ is the maximum phase speed 
of the medium. The 
pseudo-spectral method is unique among the discrete grid methods in that 
the sampling requirements are optimal, so that 
the minimum grid spacing is chosen according to ${\mit \Delta }{ x}_{min} 
\le \lambda_{min}/2 \le c_{min}/(2 \nu_{max})$ where  $\lambda_{min}$ is 
the shortest spatial  wavelength, $c_{min}$ is the minimum phase speed 
and $\nu_{max}$ is the maximum temporal frequency of the seismic source, 
consequently the time step is constrained by  
\[
{\mit \Delta}t \le 0.13 \frac{{c}_{min}}{{c}_{max} \nu_{max}}
\]
If the temporal length of the seismic record is $T$, then the number of pseudo-spectral time 
steps $N_t$ is given by
\[
N_t=\frac{T}{{\mit \Delta}t}\ge 
 \frac{T \nu_{max} {c}_{max}}{0.13 {c}_{min}}.
\]
This should be compared with our multiple scattering method, where most of 
the computation takes place in the 
{\em complex frequency domain\/} followed by an inverse discrete Fourier 
transform to synthesize 
the temporal response. The number of frequencies $N_f$ at the {\em 
Nyquist\/} sampling rate is then
\[
N_{f}=T \nu_{max}.
\]
Note that, although the frequency range is really $-\nu_{max} \le \nu \le 
\nu_{max}$, we need only compute those frequencies in the range $0 \le \nu \le 
\nu_{max}$, because the seismic response at a negative frequency is 
the complex conjugate of the response at the corresponding positive 
frequency.
Then, the number of time samples $N_{t}$ and the number of frequency 
samples $N_{f}$ are related by
\[
N_{t}=\frac{{c}_{max}}{0.13 {c}_{min}} N_{f}.
\]
For {\em gradient zones\/}, ${{c}_{max}}/{0.13 {c}_{min}}$ will be 
appreciable, so that $N_{f}\ll N_{t}$. Also, {\em some degree of  
over sampling is commonly employed in the pseudo-spectral method, 
because even though the wavefield may only require two grid points per 
wavelength, at such coarse sampling rates the propagation environment 
will be poorly resolved.\/} If a double Nyquist sampling rate is employed, 
for a medium with strong velocity gradients, then there may be $100$ 
times more time steps required than the number complex frequency points 
in the multiple scattering method. Also, the artificial means employed 
to approximate free surface and radiation conditions in the pseudo-spectral 
method requires that the computational domain be enlarged 
beyond that needed to represent the region of interest.

In {\em Marfurt\/}~(1984), the relative merits of finite difference 
algorithms for the solution of  the elastic 
wave equations in the time domain and frequency domain were compared. 
It was concluded that frequency domain (i.e. implicit) finite difference 
(FD) methods are often preferable to 
(explicit) time domain FD methods. Of course, our multiple scattering method is 
inherently a frequency domain method, indeed--multiple scattering does 
not make sense in the time domain, because causality requires that 
the scattering interactions at two distinct points in space-time are  dependent. 
Consequently, both the multiple scattering (MS) and the 
implicit FD methods require the solution of large linear 
systems at each frequency. This linear system is sparse and banded in the 
implicit FD method, but it is fully populated in the MS method. The size 
of the linear system to be solved is similar for both methods. 
Consequently, it would appear that our MS method would not be 
computationally competitive with implicit FD methods. However, this 
comparison overlooks the relative conditioning of the linear systems. We 
will take great pains to formulate the MS problem as a {\em compact linear 
integral operator equation\/}, which means that it is relatively easy to find 
numerical discretizations of arbitrarily low truncation error (Milne, 
1980). In particular, compact operators enjoy the favorable 
characteristic that a numerical representation of the operator remains 
well conditioned as the numerical mesh becomes finer. This is certainly 
not the case with the implicit finite difference approximation to the 
elastic wave equations. The elastic wave equations are differential 
operators, and consequently, they are unbounded. The unbounded nature of 
the problem is manifest by the deterioration in the conditioning of the 
numerical discretization of the differential formulation, when 
the grid size becomes finer. Hence, the linear systems that result 
from the implicit finite difference discretization of the elastic wave 
equations are more poorly conditioned than the equations that arise from 
the numerical discretization of the integral formulation of the problem.

We recall that the conditioning of a linear system is a statement about 
the complex plane distribution of the eigenvalues of that system. In 
poorly conditioned problems, the eigenvalues are spread out over the plane,
while they are bunched together in well conditioned problems (Watkins, 
1994). The distribution of the eigenvalues directly impacts on the 
efficiency of iterative schemes for the solution of linear systems. Due 
to its poor conditioning, numerical practitioners often employ {\em 
preconditioning operators\/} to the FD equations to modify its spectrum 
(Bruaset, 1995). We recall that a preconditioning operator is an 
approximate inverse of the linear system whose solution is required. 
Indeed, if we employed the exact inverse as the preconditioning operator, 
the preconditioned system would then require the solution of a linear system
with an identity matrix (all eigenvalues $\lambda$ are bunched at 
$\lambda=1$), and any consistent iterative scheme would converge in a 
single iteration. 

The great advantage of the integral (i.e. 
multiple scattering) formulation  of the elastic wave problem is that 
preconditioning is built into the equations from the outset. In the MS 
formulation, we employ a Green's tensor for the reference medium, to 
propagate the wave field along its multiple scattering path. But the 
reference Green's tensor is the {\em inverse\/} of  the elastic wave 
equations that model wave propagation in the reference medium, and 
consequently, it is a viable preconditioning operator. This is the main 
reason that the multiple scattering equations are better conditioned than 
the corresponding equations that result from an implicit finite 
difference discretization of the elastic wave equation.



However, the apparent   gains won by employing our (integral)  multiple scattering 
method should be tempered by the 
realization that  each 
successive term of the scattering series requires a single frequency, multi-source receiver 
seismogram. The number of terms in the series is usually about $20$ for a 
model with $10\%$ heterogeneity spanning the spatial region $10 \lambda_{eff} \times  10 
\lambda_{eff}$, where $\lambda_{eff}$ is the effective wavelength at the 
dominant frequency. Like the other discrete grid methods based on finite 
differences, and finite-elements, our method becomes very expensive at 
high temporal frequencies.

Calculating one frequency with our 
method, is considerably more expensive than calculating a single time 
step with the pseudo-spectral method. This is exacerbated by the more 
stringent spatial sampling requirements of our method. In particular, the lattice spacing ${\mit 
\Delta}$ of the discretized secondary source field, needed to model the scattering from 
the residual structure, must be fine on the scale $\lambda_{eff}$, say 
$\lambda_{eff}/16 \le {\mit \Delta} \le \lambda_{eff}/8$. A wavefield is 
sensitive to variations in its propagation environment at all scale 
lengths, not just the scale lengths that can be resolved at the Nyquist 
sampling rate of the wavefield. However, for practical purposes, 
we must place a resolution limit on our 
representation of the propagation environment. 


If the {\em spatial sampling rate\/} ${\mit \Delta}$ is chosen as
${\mit \Delta}=\lambda_{eff}/8$,  the number of lattice points 
$N_{\mbox{scat}}$, needed to model a region encompassing $10 \lambda_{eff} 
\times 10 \lambda_{eff}$,  is $N_{\mbox{scat}}\ge 6400$. We will soon see, that
each such scattering point has $8$ degrees of freedom, to give at 
total
of at least $N=8\times N_{\mbox{scat}}=51200$ unknowns.
 Since each scattering point, can 
interact with every other scattering point, (including itself), we get at least 
$N^2\approx 2.6 \times 10^9$ possible interactions at
 each successive multiple scattering interaction of the incident wavefield. To directly 
account for each of these interactions in turn is impractical, and this is 
certainly, the main reason that simple multiple scattering technologies
described here have not been used in computational seismology to date.

We avoid the enormous computational cost of calculating the $N^2$ 
interactions directly, by {\em factoring the wavefield\/}. We propose two 
factorizations that are complementary in their range of optimum utility:
\begin{itemize}
\item we propose an {\em alternating direction\/} sweep method
based on a Riccati factorization of the wavefield, that accounts 
for all possible interactions between points on a regular $2-$D lattice 
(over depth and offset), in less than $ O(N^2)$ floating point operations. 
Using this wavefield 
factorization, the cost of calculating the
$2.6 \times 10^9$ interactions is significantly reduced. For 
bigger problems, the reduction in computation is even larger.
{\em Riccati factorization\/} schemes for discretely layered media are well known
[{\em Chin, Hedstrom \& Thigpen\/}~(1984), {\em Ha\/}~(1985)], and are 
implicit in the {\em invariant embedding schemes\/} of {\em 
Kennett\/}~(1984), {\em Tromp \& Snieder\/}~(1988).
\item  a {\em cumulative forward 
scattering \/} factorization, that generalizes the method of {\em Wu\/}~(1994) 
to a weakly stratified reference halfspace. In this method we march the forward 
scattered field.
\end{itemize}
The Riccati sweep method, is a two way marching scheme that factorizes 
the reference medium Green's tensor, but does not include any of the 
scattered field in the reference solution. Hence, a multiple scattering 
series based on this scheme converges rapidly, even when there are strong 
gradients in the reference model, provided that the superimposed 
heterogeneity is not too strong. The complementary problem, where the 
reference model is rather smooth, but the superimposed heterogeneity may 
be strong is best handled by the cumulative forward 
scattering factorization, as the forward scattered energy is incorporated 
directly into the reference solution, and we sum a multiple scattering 
series for the neglected back scattered field.


A {\em fast multipole method\/} has recently been proposed 
[{\em Rokhlin\/} ~(1985), {\em Rokhlin\/} ~(1990)] that is based on {\em 
classical potential theory\/} for the {\em Laplace equation\/}. This 
method can compute an approximation to the $N^2$ gravitational 
interactions between $N$ point masses, in very much less than $N^2$ 
floating point operations. In the 
problem addressed in this thesis, we generate the wave field by 
spatial differentiation of  solutions of the {\em Helmholtz equation\/}, 
so that much of the theory proposed in {\em Rokhlin\/}~(1985) and {\em 
Rokhlin\/}~(1990) can be adapted to our problem. The main problem with the existing fast 
multipole theory, is that it is based on the {\em spatial domain free space
Green's function\/}, so that it is not readily extensible to plane 
stratified media. A second, but  less severe problem, is the very sophisticated 
mathematics, and computer programming needed to implement the method.


The main advantages of the multiple scattering approach adopted here, over competing 
technologies, based directly on discretization of the elastic wave
equations,[e.g. {\em Kosloff et al.\/}~(1984), {\em Reshef \& 
Kosloff\/}~(1985), {\em Vireaux}~\shortcite{Vir86}, {\em 
Levander}~\shortcite{Lev88}] are:
\begin{itemize}
\item {\em free-surface, and radiation conditions\/} are 
built into the development from the outset

\item it is straightforward to introduce a point source

\item attenuation (both causal and acausal) are readily accommodated

\item the components of the model have a direct and clear interpretation
in terms of the  physics of the underlying wave propagation processes, 
allowing the interpreter great flexibility in his/her choice of model
e.g. it is straightforward to 
isolate single scattered arrivals from multi-scattered arrivals, or separate
body wave from surface waves, etc.

\item the numerical discretization of our problem has spectral properties that are 
favourable for numerical solution via iterative methods, whereas 
discretizations based directly on the differential formulation are 
notoriously difficult to solve using iterative solution methods
\end{itemize}
 

%%%%%%%%%%%%%%%%%%%%%%%%%%%%%%%%%%%%%%%%%%%%%%%%%%%%%%%%%%%%%%%%%%%%%%%%%%%%%%%%%%%%%%%%%%%%
%%%%%%%%%%%%%%%%%%%%%%%%%%%%%%%%%%%%%%%%%%%%%%%%%%%%%%%%%%%%%%%%%%%%%%%%%%%%%%%%%%%%%%%%%%%%
%%%%%%%%%%%%%%%%%%%%%%%%%%%%%%%%%%%%%%%%%%%%%%%%%%%%%%%%%%%%%%%%%%%%%%%%%%%%%%%%%%%%%%%%%%%%

\section{the basic formalism}
Our method is but one of a whole class of methods, that attempt to construct the solution of the 
given elastic wave equations in terms of a soluble allied problem. This method is 
rather flexible, as it applies to the entire class of {\em linear 
field equations\/}, although it assumes a multitude of manifestations, 
depending upon the precise nature of the relationship between the given 
problem and its allied counterpart. The success of the method hinges,
 on the ability solve an allied problem, whose solution 
reflects the actual physics of the given problem.

A recent application of this concept in the context of high 
frequency body wave propagation was proposed by {\em  Coates \& Chapman\/} 
~(1991), where the allied problem was chosen as the high frequency 
asymptotic approximation of the elastic wave equation. These authors went 
on to develop a {\em representation theorem\/}, that expressed the elastic 
wavefield in terms of the solution of the asymptotic allied problem. The 
advantage of their approach is that the allied solution can be expected 
to be a fairly good approximation to the required elastic wavefield, at 
least for high frequencies and smooth variations of material properties,
 so that the equations that result from their 
representation theorem need only be solved to within the Born approximation.
A potential disadvantage is that the equation error of the 
asymptotic allied problem, does not have compact support, and this 
complicates the representation theorem through the need to incorporate 
complicated boundary integrals. Also, it would appear rather difficult to 
solve the integral equations beyond the level of the first Born 
approximation.

For the bulk of this thesis we follow the more traditional route, where the allied 
problem is chosen as the elastic wave equation, but assume that the 
material variations are simple enough that the solution of the allied 
equation may be explicitly constructed. The added material variations, 
not already included in the reference model, then act as a {\em secondary 
or virtual source field\/} to be superimposed on the reference field. As 
pointed out by {\em Coates \& Charrette\/}~(1993), this approach has the 
significant disadvantage that the incident field at each scattering 
interaction is {\em habitually mis-timed\/}. For the most part, we have chosen to 
live with this deficiency and focus attention instead,   on the development of 
{\em iterative schemes\/} that are robust enough to converge the multiple 
scattering series representation of the secondary source field. This task 
is aided by the development of Riccati factorization schemes that 
allow us the propagate the wavefield between successive scattering 
interactions in an efficient and stable manner. 

A recent publication (Wu, 1994), 
concerning {\em the elastic wave phase screen method\/} has addressed 
some of the very issues of concern in our method.  The method of {\em Wu\/}~(1994) is 
basically a single scattering method, where the {\em cumulative forward 
scattered field\/} is marched forward. In the most basic version of our 
multiple scattering  
method, we sum the multiple 
scattering series as a Born series, which is equivalent to 
solving the matrix representation of the scattered field using a {\em Jacobi 
iteration\/} scheme. Instead, {\em Wu\/}~(1994) employs 
what is equivalent, to a single iteration of the {\em Gauss-Seidel\/} scheme. 
In the Gauss-Seidel version,  the forward scattered reference field  has the 
correct phase for direct interaction paths between successive scattering 
points, which corrects for the {\em habitual mis-timing\/} in our basic 
(Born series) method. This is at the cost of employing a reference 
solution that ignores both the reflected 
field in the reference model, as well as the backscattered field from the 
residual model.
The consequence of this alternate choice of reference field is that the 
multiple scattering series for the secondary source field converges much faster, 
than the corresponding series for 
a reference model that does not include any forward scattered energy, in 
those problems where reflected phases in the reference model are 
unimportant (i.e. in smooth reference models).
We employ a set of {\em wavefield balances\/} to find a representation of 
the backscattered field, in so doing,
we  have extended the method of {\em Wu\/}~(1994) 
into a complete wave theory, in the same way that {\em Coates \& 
Chapman\/}~(1991), extend ray theory. 

However, for the bulk of the analysis and computation, we have 
concentrated on the two-way Riccati factorization approach, where the 
reference Green's tensor is the exact solution for a plane stratified 
halfspace, and {\em none\/} of the scattered field has been incorporated 
within it. The burden of convergence of the multiple scattering series is 
then borne by the method used to sum the multiple scattering partial sums. 
Herein, lies the principal contribution of this research. We have 
employed two summation schemes, that exploit complementary properties of 
the wavefield. 

The first summation method is a generalization of the {\em 
Schwinger-Levine\/} variational method introduced into seismology by {\em 
Lapwood \& Hudson\/}~(1975). A Schwinger-Levine variational principle is 
an expression of the symmetry of a scattered wavefield when the source point and 
the field point are interchanged, i.e. {\em a statement of scattered 
wavefield reciprocity\/}. 
Mathematically, the primal and reciprocal scattering experiments are {\em 
dual\/} statements of the same physical problem. More explicitly,  if we introduce a 
transition operator ${\CB T}$, then the {\em scattering transition\/} 
from the reference medium  incident state  ${\bmi b}_{inc}$ to 
reference medium the scattered state  ${\bmi 
b}_{sc}$ is modelled as
\[
\left\langle{\bmi b}_{sc}, {\CB T}{\bmi b}_{inc}\right\rangle
\]
while the {\em  reciprocal transition\/} from ${\bmi b}_{sc}$ to 
 ${\bmi 
b}_{inc}$ is modelled as
\[
\left\langle {\CB T}^{\dagger}{\bmi b}_{sc}, {\bmi b}_{inc}\right\rangle
\]
where ${\CB T}^{\dagger}$ is the {\em adjoint transition operator\/} 
relative to the bilinear form $\langle \cdot, \cdot \rangle$ that 
expresses the scattering transitions induced by the lateral heterogeneities
superimposed on the reference medium.

In {\em Lapwood \& Hudson\/}~(1975),  a bilinear form was 
introduced for the plane stratified problem. We have generalized this 
bilinear form to the heterogeneous problem. The second extension of  
research of {\em Lapwood \& Hudson\/}~(1975)
concerns extending the accuracy of their method beyond  second 
order Born theory. The key step in this extension is to regard their 
variational expression as a low order {\em Pad\'{e}} approximation to the 
true multiple scattering series representation of the scattered field. The 
higher order corrections to the leading order expression of {\em Lapwood \& Hudson\/}~(1975)
may then be derived by writing down the relevant higher order {\em Pad\'{e}} approximation to the 
true multiple scattering series. More importantly, these high order {\em 
Pad\'{e}} approximations can be calculated by exploiting their well known 
relationship to the {\em bi-conjugate gradient method\/}. Using this line 
of reasoning, we have extended the {\em Lapwood \& Hudson\/} variational 
method to arbitrarily high orders of perturbation theory. However, this 
method is deficient in two respects:
\begin{itemize}
\item due to the dependence of the method on seismic reciprocity, each 
source/receiver combination requires the solution of a separate  problem
\item the bilinear form that expresses the scattering transition is 
indefinite. Consequently, convergence of the method is not guaranteed
\end{itemize}

To avoid these deficiencies, we have also generalized the {\em 
Generalized Minimal Residual Method\/}~(GMRES) of {\em Saad \& 
Schultz\/}~(1986) which is based on solving a set of {\em linear least-squares 
minimization problems\/}, optimized over a trial space spanned by the multiple 
scattering series partial sums. In each iteration of this method, the 
trial space is augmented by an extra multiple scattering interaction. The 
advantage of the  GMRES method is that the residual of the current solution may not increase 
from one iteration to the next, so that the convergence of the method is 
very smooth. Also, the method calculates the response at all 
receivers from a given source point, simultaneously. Our generalization 
of the basic method involves:
\begin{itemize}
\item the extension to complex arithmetic
\item the use of a recursive factorization scheme based on {\em unitary 
Householder reflections\/}, to solve the linear least squares problem at 
each iteration of the GMRES method
\end{itemize}

%%%%%%%%%%%%%%%%%%%%%%%%%%%%%%%%%%%%%%%%%%%%%%%%%%%%%%%%%%%%%%%%%%%%%%%%%%%%%%%%%%%%%%%%%%%%%
%%%%%%%%%%%%%%%%%%%%%%%%%%%%%%%%%%%%%%%%%%%%%%%%%%%%%%%%%%%%%%%%%%%%%%%%%%%%%%%%%%%%%%%%%%%%%
%%%%%%%%%%%%%%%%%%%%%%%%%%%%%%%%%%%%%%%%%%%%%%%%%%%%%%%%%%%%%%%%%%%%%%%%%%%%%%%%%%%%%%%%%%%%%


\section{notation}
In this thesis, we make use of both the {\em Kronecka product\/} between 
matrices, and the {\em tensor product\/} between vectors. Although both algebraic 
operators are represented by the binary operator ``$\otimes$", they 
should considered as distinct mathematical operations.
The  Kronecka product ${\bmi A} \otimes {\bmi B}$  of an $m \times n$ matrix $\bmi A$ and a 
$p \times q$  matrix ${\bmi B}$ is the $(m+p) \times (n+q)$ matrix
\[
{\bmi A} \otimes {\bmi B} \equiv \left[{\begin{array}{ccc}{ 
A}_{ 11}{\bmi  B}&\cdots &{ A}_{1n}{\bmi B}\\ \vdots&\cdots &\vdots\\ 
{A}_{m1}{\bmi B}&\cdots &{A}_{mn}{\bmi B}\end{array}}\right]
\] 
whereas the {\em tensor\/} product ${\bmi a} \otimes {\bmi b}$
of an $m-$vector $\bmi a$ and a 
$p-$vector ${\bmi b}$ is the $m \times p$ matrix
\[
{\bmi a} \otimes {\bmi b} \equiv \left[{\begin{array}{c}{ 
a}_{ 1}{\bmi  b}^{T}\\ \vdots\\ 
{a}_{m}{\bmi b}^{T}\end{array}}\right]
\] 
so that {\em there is an implied transposition in the tensor product between 
vectors, but there is no implied transposition in the Kronecka product 
between matrices\/}. The mixed cases ${\bmi A} \otimes {\bmi b}$ and  
${\bmi b} \otimes {\bmi A} $ will always be interpreted as {\em 
Kronecka\/} products, so there is no implied transposition in these cases.


